% !TEX program = xelatex

\documentclass{resume}

% optionally suppress printing the page number
\pagenumbering{gobble}

% optional Chinese support
% \useChinese{Medium}{Heavy}
\useChinese{Regular}{Bold}
\setstretch{1.12}
\titlespacing*{\section}{0pt}{0.45em}{0.45em}
\newcommand{\projectgap}{\vspace{0.55em}}
\newcommand{\itemsummary}[1]{{\setlength{\parindent}{12pt}\small\textit{\textcolor{gray}{#1}}\par}}

\begin{document}

%% basic personal info
\name{\textbf{WANG, Yankai\hspace{0.4em}\large（ワン・イエンカイ）}}
\begin{basicinfo}
	\info{\location{Tokyo, Japan}}
	\info{\email{ykwang224@gmail.com}}
	\info{\homepage{shiinayane.com}[https://shiinayane.com]}
	\info{\github{shiinayane}[https://github.com/shiinayane]}
\end{basicinfo}

\projectgap
\noindent Japanese (Professional working, JLPT N1 121/180) \, $\cdot$ \, English (Professional, TOEFL 97/120) \, $\cdot$ \, Chinese (Native)

\section{\texorpdfstring{\faGraduationCap\ }{}\textbf{Education}}

\cventry{The University of Tokyo\hspace{0.4em} \normalsize 東京大学}
{Oct. 2025 -- Sep. 2027 (Expected)}
[\textit{M.S. Candidate}, Information \& Communication Engineering]
[Tokyo, Japan]
\begin{itemize}
	\item Suzumura Laboratory; pursuing research on \textbf{Vision-Language-Action (VLA)} models and human-robot collaboration benchmarks for scientific laboratory environments.
\end{itemize}

\cventry{Shanghai University\hspace{0.4em} \normalsize 上海大学}
{Sep. 2021 -- Jul. 2025}
[\textit{B.E.}, Electronic Information Engineering]
[Shanghai, China]

\section{\texorpdfstring{\faProjectDiagram\ }{}\textbf{Projects}}

\cventry{KotobaLab}
{Apr. 2026 -- Present}
[\textbf{Personal Project}: SwiftUI, GRDB 7, SQLite, SwiftData, Swift Testing, iOS 18+]
[Tokyo, Japan]
\itemsummary{Local-first Japanese dictionary app built with SwiftUI, MVVM, and a Clean Architecture–inspired layered design.}
\itemsummary{Currently in local development.}
\begin{itemize}
	\item Built a local-first dictionary app with a \textbf{293K-entry, 52\,MB SQLite} database generated by a Python pipeline, bundled as an app resource and queried through GRDB \texttt{DatabaseQueue}.
	\item Optimized prefix search with \texttt{PRAGMA case\_sensitive\_like=ON}, improving benchmarked queries from \textasciitilde16\,ms table scans to \textbf{\textasciitilde0.03\,ms} multi-index lookups.
	\item Designed a four-layer architecture (\texttt{AppDependencies} $\rightarrow$ \texttt{Scene} $\rightarrow$ \texttt{Store} $\rightarrow$ \texttt{View}) following \textbf{MVVM} and \textbf{Clean Architecture} principles; applied the \textbf{Repository pattern} with protocol-based abstractions and dependency injection to decouple Domain logic from UI and storage frameworks.
	\item Combined GRDB-managed SQLite for read-heavy reference data with SwiftData for user state; wrote 13 unit tests with Swift Testing covering domain use cases and repository behavior.
	\item Roadmap: SQLite FTS5 full-text search, actor-isolated repositories under Swift Concurrency, TestFlight beta.
	\item Github: \iconlink[\faGithub][shiinayane/KotobaLab]{https://github.com/shiinayane/KotobaLab}.
\end{itemize}

\projectgap

\cventry{LabMate}
{Dec. 2025 -- Present}
[\textbf{Suzumura Lab}: Python, Isaac Sim, VLA Models]
[Tokyo, Japan]
\itemsummary{Simulation-first benchmark for human-robot collaboration in scientific laboratory environments.}
\begin{itemize}
	\item Extending the \textbf{LabUtopia} simulator into a benchmark for natural-language human-robot collaboration in scientific laboratories, focusing on clarification, safety-aware refusal, grounding, and structured action logging.
	\item Designing evaluation framework to compare planner architectures (\textbf{LLM}, \textbf{VLA}, and hybrid systems) on multi-step laboratory protocols; sim-first phase planned for 6+ months before hardware integration.
	\item Coordinating with 3 procured robotic platforms for future hardware integration; prioritizing rigorous evaluation methodology over demonstration-driven research.
\end{itemize}

\section{\texorpdfstring{\faUsers\ }{}\textbf{Experience}}

\cventry{GlucoPI}
{Feb. 2025 -- Jun. 2025}
[\textbf{B.E. Graduation Project}, \textbf{Advisor}: \href{https://scie.shu.edu.cn/Prof/zhangq.htm}{\textit{Qi Zhang}}, PhD, Professor]
[Shanghai, China]
\itemsummary{WeChat Mini Program for diabetes management, doctor-patient communication, and glucose prediction.}
\begin{itemize}
	\item Built an end-to-end diabetes self-tracking app with WeChat Mini Program frontend, FastAPI backend, MySQL/MongoDB storage, and WebSocket doctor-patient chat.
	\item Implemented login and permissions, glucose / diet / insulin records, doctor-patient pairing, real-time messaging, and trend charts.
	\item Added short-horizon glucose prediction in PyTorch by reproducing the GluPred pipeline and running experiments on \textbf{OhioT1DM}.
	\item Integrated a lightweight LLM Q\&A feature for non-diagnostic health guidance.
	\item Github:\iconlink[\faGithub][shiinayane/glucopi]{https://github.com/shiinayane/glucopi}.
\end{itemize}

\projectgap

\cventry{Arrived or Not}
{Dec. 2023 -- Dec. 2024}
[\textbf{Role}: Co-Leader, \textbf{Advisor}: \href{https://scie.shu.edu.cn/Prof/zhangq.htm}{\textit{Qi Zhang}}, PhD, Professor]
[Shanghai, China]
\itemsummary{Machine-vision teaching platform for classroom attendance and course management.}
\begin{itemize}
	\item Implemented real-time face-recognition attendance with RetinaFace detection and \textbf{ArcFace} embedding-based identity matching.
	\item Built teacher / student mobile workflows for class setup, attendance records, and basic engagement statistics.
	\item Developed a Flutter client and FastAPI service with a clear frontend-backend split.
	\item Github: \iconlink[\faGithub][shiinayane/Arrived-or-Not-Frontend]{https://github.com/shiinayane/Arrived-or-Not-Frontend} and \iconlink[\faGithub][shiinayane/Arrived-or-Not-Backend]{https://github.com/shiinayane/Arrived-or-Not-Backend}.
\end{itemize}

\section{\texorpdfstring{\faLightbulb\ }{}\textbf{Interests}}

\begin{compactlist}
	\item \textbf{Product-First Engineering}: Building products that solve real problems; treating iOS, backend, and full-stack as means rather than ends.
	\item \textbf{iOS Native Architecture}: Local-first design with clean Domain / UI separation; currently focused on shipping native iOS apps as the most direct path to product impact.
	\item \textbf{Research-to-Product}: Bridging VLA / LLM research from graduate work to consumer mobile experiences.
\end{compactlist}

\section{\texorpdfstring{\faCogs\ }{}\textbf{Skills}}

\begin{compactlist}
	\item \textbf{iOS / Apple Platforms}: Swift, SwiftUI, GRDB, SQLite, SwiftData, Observation, Swift Concurrency, Swift Testing
	\item \textbf{Architecture \& Design}: MVVM, Clean Architecture, Repository Pattern, Dependency Injection, Protocol-Oriented Programming
	\item \textbf{Backend / Web}: Python (FastAPI), JavaScript
	\item \textbf{ML / Research}: PyTorch, VLA Models, Isaac Sim
	\item \textbf{Tools / Platforms}: Git, Docker, Xcode, VS Code, macOS, Linux
\end{compactlist}

\end{document}
